
\documentclass[12pt,a4paper]{article}

\usepackage[utf8]{inputenc}
\usepackage[T1]{fontenc}
\usepackage[spanish]{babel}
\usepackage[margin=2.5cm]{geometry}
\usepackage{enumitem}
\usepackage{hyperref}

\title{FORM-TL-01\\
Notas de presentación: Vecindario}

\author{
APUNTE PAZMIÑO WASHINGTON GERMAN\\
LOOR SUAREZ JORDAN JOSUE\\
MACIAS VASQUEZ CESAR ABRAHAN
}

\date{Periodo académico 2026--2027}

\begin{document}

\maketitle

\section{Contexto del escenario}

Vecindario es una red social ficticia orientada a comunidades
de barrio en ciudades intermedias de Ecuador.

La plataforma permite publicar avisos, seguir a vecinos y
negocios locales, reaccionar, comentar y mantener conversaciones
mediante mensajes directos.

\section{Arquitectura propuesta}

La arquitectura propuesta para Vecindario se organiza mediante servicios independientes que se encargan de las diferentes funciones de la red social. Esta separación permite distribuir el procesamiento de las publicaciones, la mensajería, las notificaciones y la moderación.

Los componentes principales son:

\begin{itemize}
    \item \textbf{Cliente (CLI):} aplicación móvil o web desde la que los usuarios acceden a la plataforma, publican avisos y utilizan sus funcionalidades.

    \item \textbf{Servicio de publicaciones (PUB):} permite crear, editar y eliminar avisos, incluyendo publicaciones con imágenes.

    \item \textbf{Servicio de seguidores (SEG):} administra las relaciones entre usuarios, negocios locales y sus seguidores.

    \item \textbf{Servicio de muro (MUR):} organiza y distribuye las publicaciones que aparecen en el muro de cada usuario.

    \item \textbf{Servicio de mensajería directa (MSG):} administra las conversaciones entre usuarios y permite conocer su estado de conexión.

    \item \textbf{Servicio de moderación (MOD):} recibe las denuncias y gestiona el ocultamiento de publicaciones que incumplen las normas.

    \item \textbf{Servicio de notificaciones (NOT):} comunica a los usuarios los avisos y eventos relevantes de la plataforma.
\end{itemize}

Además, intervienen los siguientes actores externos:

\begin{itemize}
    \item \textbf{Usuario:} utiliza la red social para publicar avisos, consultar el muro e intercambiar mensajes.
    \item \textbf{Negocio local:} publica información de sus productos, servicios y actividades para llegar a sus seguidores.
    \item \textbf{Moderador:} revisa las denuncias y participa en el proceso de moderación.
    \item \textbf{Servicio externo de almacenamiento de imágenes:} conserva las imágenes utilizadas en las publicaciones.
\end{itemize}

La arquitectura también contempla componentes de infraestructura para gestionar los eventos, conservar los mensajes, consultar la presencia de los usuarios y distribuir las conexiones de mensajería.

\section{Mecanismos de comunicación}

% Documentar sockets, RPC y mensajería.

\section{Protocolos y middleware}

% Registrar los protocolos y middleware seleccionados,
% junto con su justificación técnica.

\section{Requisitos y situaciones de fallo}

% Explicar cómo el diseño responde a R1-R6 y F1-F3.

\section{Preguntas críticas recibidas y respuestas}

\begin{enumerate}[leftmargin=*]

    \item \textbf{Pregunta recibida:}

    \textbf{Equipo que la formuló:}

    \textbf{Respuesta del equipo:}

\end{enumerate}

\section{Preguntas críticas formuladas a otros equipos}

\begin{enumerate}[leftmargin=*]

    \item \textbf{Equipo al que se dirigió:}

    \textbf{Pregunta formulada:}

\end{enumerate}

\section{Retroalimentación y ajustes}

% Registrar las observaciones del docente y los cambios
% realizados al diseño después de la presentación.

\end{document}
